\section*{Results}

\subsection*{Delivery performance in and out of distribution}

Following the protocol described in the Methods, every policy was trained on the nominal distribution with three seeds and evaluated on paired scenarios: each policy meets the same district layouts, customers, wind fields, battery levels and payloads, so differences come from behaviour rather than from luck of the draw. Table~\ref{tab:main} reports the six test conditions with the conventional fixed reserve, and Fig.~\ref{fig:main} shows the catastrophic loss rate, the fraction of missions ending in a collision, an energy depletion or an airspace exit, for both reserve rules.

Training curves on the nominal distribution are shown in Fig.~\ref{fig:learn}. On the nominal distribution the five off-policy agents are close together. The soft actor-critic loses \RSacStdLossNominal\% of missions, the risk-neutral distributional agent \RDsacnStdLossNominal\%, the fixed-risk CVaR agent \RDsaccvarStdLossNominal\%, TD3 \RTdthreeStdLossNominal\% and CREST-C \RCrestcStdLossNominal\%, with delivery success between \RTdthreeStdSuccNominal\% and \RDsacnStdSuccNominal\%. On-policy PPO is far behind at this budget, with a loss rate of \RPpoStdLossNominal\% and success of \RPpoStdSuccNominal\%, which reflects the sample cost of on-policy learning in a multi-agent setting rather than a defect of the algorithm.

All four shift families degrade every policy, and they do so through different mechanisms. Under the wind shift the loss of the soft actor-critic rises to \RSacStdLossWind\%, of which \RSacStdDepWind\ percentage points are energy depletion rather than collision. Under the sensor shift, in which navigation error grows to five metres with frequent multipath excursions, building strikes rise from \RSacStdCollbNominal\% to \RSacStdCollbSensor\%. Under the traffic shift, which doubles the fleet and more than doubles the intruder count, collisions with other aircraft dominate, reaching \RSacStdCollaTraffic\% for the soft actor-critic and \RCrestcStdCollaTraffic\% for CREST-C. Under the energy shift, which lowers battery capacity and raises payload and fault rates, depletion becomes the leading failure for the fixed reserve at \RSacStdDepEnergy\%, and delivery success falls to \RSacStdSuccEnergy\%. The compound condition, which applies all four families at half severity, is the hardest overall.

\subsection*{Tail-risk objectives give no reliable reduction in catastrophic losses}

The central comparison of this study is between the risk-neutral agents and the three tail-risk designs at a matched training budget. The differences are small and inconsistent in sign. With the fixed reserve, the catastrophic loss rate of CREST-C is \RCrestcStdLossNominal\% at nominal against \RSacStdLossNominal\% for the soft actor-critic, \RCrestcStdLossSensor\% against \RSacStdLossSensor\% under the sensor shift, \RCrestcStdLossTraffic\% against \RSacStdLossTraffic\% under the traffic shift and \RCrestcStdLossCompound\% against \RSacStdLossCompound\% under the compound shift. The corresponding bootstrap intervals overlap in every condition; for example the nominal interval is [\CCrestcStdLossNominalLo, \CCrestcStdLossNominalHi] for CREST-C and [\CSacStdLossNominalLo, \CSacStdLossNominalHi] for the soft actor-critic. The fixed-risk CVaR agent behaves like its risk-neutral counterpart, with \RDsaccvarStdLossNominal\% against \RDsacnStdLossNominal\% at nominal and \RDsaccvarStdLossTraffic\% against \RDsacnStdLossTraffic\% under the traffic shift. Across seeds the standard deviation of the loss rate is of the same order as the gaps between methods, for instance \SSacStdLossNominal\ percentage points for the soft actor-critic and \SCrestcStdLossNominal\ for CREST-C at nominal.

The risk level can be moved at deployment time because the policy is risk-conditioned, and Fig.~\ref{fig:dial} with Table~\ref{tab:dial} sweep it. Lowering the base risk level from \(\alpha_b=1\) to \(\alpha_b=0.1\) does not reduce collisions: under the sensor shift the collision rate moves from \DOnepZeroCollSensor\% to \DZeropOneCollSensor\%, and under the traffic shift from \DOnepZeroCollTraffic\% to \DZeropOneCollTraffic\%. Delivery time is nearly flat across the dial, changing from \DOnepZeroTdelNominal\~s to \DZeropOneTdelNominal\~s at nominal. Figure~\ref{fig:alpha} shows the schedule doing what it was designed to do along a single delivery, tightening the risk level near conflicts and as the energy margin falls. Inspection of trajectories shows why the behavioural effect is not the intended one: a lower risk level makes the aircraft hesitant, reducing speed and stopping short of conflicts, which lengthens the time spent near other traffic instead of widening separation. Because the exposure grows while the geometric margin does not, the tail of the outcome distribution is unchanged.

The ablations in Fig.~\ref{fig:abl} and Table~\ref{tab:abl} isolate the components of the risk regulation. Replacing the hazard-adaptive schedule with a constant risk level changes the loss rate by less than a percentage point in most conditions (\AFixedalphaLossNominal\% against \RCrestcOnlineLossNominal\% at nominal and \AFixedalphaLossTraffic\% against \RCrestcOnlineLossTraffic\% under the traffic shift). Removing the energy cue, keeping only proximity, or replacing the cue-driven schedule by one driven by the spread of the learned return distribution, which is the mechanism used by adaptive risk-tendency methods, all produce similarly small changes. The same holds for the variant that places the tail-risk measure on the return distribution rather than on the catastrophe cost.

\subsection*{Shift severity and the composition of failures}

Figure~\ref{fig:sev} and Table~\ref{tab:sev} sweep each shift family from nominal to full severity. The curves are ordered by algorithm family rather than by risk sensitivity: the off-policy agents track each other within a few percentage points throughout, and PPO sits far above them. The energy family produces the steepest degradation in delivery success, the traffic family the steepest rise in collisions, and the sensor family degrades building clearance without affecting energy. Nothing in the sweeps suggests that risk-sensitive training buys robustness that accumulates with severity. Efficiency differences between the off-policy agents are equally modest: mean energy per mission spans \RDsacnStdEnergyNominal--\RSacStdEnergyNominal~kJ at nominal and \RDsacnStdEnergyCompound--\RSacStdEnergyCompound~kJ in the compound condition, and mean delivery time spans \RDsacnStdTdelNominal--\RSacStdTdelNominal~s at nominal.

Figure~\ref{fig:extra} decomposes what remains. The composition of failures also explains why the risk objectives have little leverage. Collisions with other aircraft account for \RSacStdCollaTraffic\ of \RSacStdCollTraffic\ percentage points of collisions under the traffic shift, and these involve non-cooperative intruders that are detected late, or not at all, given a miss probability of one in five under the sensor shift. A policy cannot avoid what it has not observed, and the tail of its own return distribution carries no information about an undetected aircraft.

\subsection*{Calibrated reserves control energy risk}

The picture is different for the return-to-home decision. Figure~\ref{fig:cov}a shows the fraction of return legs whose energy stayed within the reserve predicted at the start of the leg. The fixed \(1.2\times\) margin covers \RCrestcStdCovNominal\% of legs at nominal and only \RCrestcStdCovEnergy\% under the energy shift, and the \(1.5\times\) margin degrades in the same way, from \RCrestcStdfifteenCovNominal\% to \RCrestcStdfifteenCovEnergy\%. The online conformal reserve holds its target in every condition, between \RCrestcOnlineCovNominal\% and \RCrestcOnlineCovTraffic\%, which is the behaviour Proposition~\ref{prop:coverage} predicts for an arbitrary sequence of flights. Figure~\ref{fig:lam} shows the fleet calibrator adapting within a few hundred legs and the running miscoverage settling at the target.

Calibration translates into fewer energy losses. With the fixed \(1.2\times\) margin, depletion under the energy shift costs \RSacStdDepEnergy\% of missions for the soft actor-critic and \RCrestcStdDepEnergy\% for CREST-C; with the online conformal reserve these fall to \RSacOnlineDepEnergy\% and \RCrestcOnlineDepEnergy\%. The cost is paid in abandoned deliveries, which rise from \RSacStdRthEnergy\% to \RSacOnlineRthEnergy\% for the soft actor-critic, so the reserve is a trade-off rather than a free improvement.

Table~\ref{tab:online} repeats the outcome table under the online conformal reserve, and Table~\ref{tab:reserve} shows the same pattern for both policies and adds the price of the reserve in energy terms: the online rule reserves about \RCrestcOnlineTightEnergy\ times the energy a return leg actually consumes under the energy shift, against \RCrestcStdTightEnergy\ for the fixed \(1.2\times\) margin, which is what converts unusable margin into abandoned deliveries. Figure~\ref{fig:pareto} and Table~\ref{tab:par} make that trade-off explicit by sweeping fixed margins from \(1.0\times\) to \(2.5\times\) against conformal reserves with target miscoverage from 2\% to 30\%. At operating points with negligible depletion the conformal reserve is the better instrument: under the energy shift it abandons \PsaconlineZeropZeroFiveRthEnergy\% of deliveries for \PsaconlineZeropZeroFiveDepEnergy\% depletion, where the fixed margin that reaches comparable safety, \(2.0\times\), abandons \PsacfixedTwopZeroRthEnergy\% for \PsacfixedTwopZeroDepEnergy\%. At nominal the same conformal setting abandons \PsaconlineZeropZeroFiveRthNominal\% against \PsacfixedTwopZeroRthNominal\% for the equally safe fixed margin. The advantage is not universal. Under the wind shift a \(1.5\times\) margin abandons \PsacfixedOnepFiveRthWind\% of deliveries at \PsacfixedOnepFiveDepWind\% depletion, which the conformal settings do not improve on, because the wind shift inflates the modelled return energy itself rather than the discrepancy the calibrator tracks. The practical reading is that a single fixed margin cannot serve all conditions, the margin that is right at nominal being unsafe under the energy shift and the margin that is safe under the energy shift being wasteful at nominal, whereas one conformal target holds its risk level everywhere at a cost that the operator can set.
